% Options for packages loaded elsewhere
\PassOptionsToPackage{unicode}{hyperref}
\PassOptionsToPackage{hyphens}{url}
\PassOptionsToPackage{dvipsnames,svgnames,x11names}{xcolor}
\documentclass[
  11pt,
]{article}
\usepackage{xcolor}
\usepackage[a4paper,margin=2.2cm]{geometry}
\usepackage{amsmath,amssymb}
\setcounter{secnumdepth}{-\maxdimen} % remove section numbering
\usepackage{iftex}
\ifPDFTeX
  \usepackage[T1]{fontenc}
  \usepackage[utf8]{inputenc}
  \usepackage{textcomp} % provide euro and other symbols
\else % if luatex or xetex
  \usepackage{unicode-math} % this also loads fontspec
  \defaultfontfeatures{Scale=MatchLowercase}
  \defaultfontfeatures[\rmfamily]{Ligatures=TeX,Scale=1}
\fi
\usepackage{lmodern}
\ifPDFTeX\else
  % xetex/luatex font selection
\fi
% Use upquote if available, for straight quotes in verbatim environments
\IfFileExists{upquote.sty}{\usepackage{upquote}}{}
\IfFileExists{microtype.sty}{% use microtype if available
  \usepackage[]{microtype}
  \UseMicrotypeSet[protrusion]{basicmath} % disable protrusion for tt fonts
}{}
\makeatletter
\@ifundefined{KOMAClassName}{% if non-KOMA class
  \IfFileExists{parskip.sty}{%
    \usepackage{parskip}
  }{% else
    \setlength{\parindent}{0pt}
    \setlength{\parskip}{6pt plus 2pt minus 1pt}}
}{% if KOMA class
  \KOMAoptions{parskip=half}}
\makeatother
\usepackage{longtable,booktabs,array}
\newcounter{none} % for unnumbered tables
\usepackage{calc} % for calculating minipage widths
% Correct order of tables after \paragraph or \subparagraph
\usepackage{etoolbox}
\makeatletter
\patchcmd\longtable{\par}{\if@noskipsec\mbox{}\fi\par}{}{}
\makeatother
% Allow footnotes in longtable head/foot
\IfFileExists{footnotehyper.sty}{\usepackage{footnotehyper}}{\usepackage{footnote}}
\makesavenoteenv{longtable}
\setlength{\emergencystretch}{3em} % prevent overfull lines
\providecommand{\tightlist}{%
  \setlength{\itemsep}{0pt}\setlength{\parskip}{0pt}}
% Preamble for Complete Notes PDF build (used by build_notes.sh via pandoc + xelatex)
% XeLaTeX is required: the notes contain umlauts, arrows and symbols.

\usepackage{fontspec}
\usepackage{xcolor}
\usepackage{booktabs}
\usepackage{longtable}
\usepackage{array}
\usepackage{fancyhdr}
\usepackage{titlesec}
\usepackage{enumitem}

% --- Fonts -------------------------------------------------------------
% Main text: a serif face with full German coverage.
\setmainfont{Palatino}
\setsansfont{Helvetica Neue}
\setmonofont[Scale=0.9]{Menlo}

% Fallback for symbols the main font lacks (arrows, check marks).
% Emoji are stripped by build_notes.sh before pandoc runs, so no emoji font
% is declared here — declaring one that may be absent breaks the build.
\newfontfamily\symbolfont{Apple Symbols}

% --- Colours -----------------------------------------------------------
\definecolor{headingblue}{HTML}{1F4E79}
\definecolor{warnred}{HTML}{B03A2E}
\definecolor{rulegrey}{HTML}{BBBBBB}

% --- Headings ----------------------------------------------------------
\titleformat{\section}{\Large\bfseries\color{headingblue}}{}{0em}{}
\titleformat{\subsection}{\large\bfseries\color{headingblue}}{}{0em}{}
\titleformat{\subsubsection}{\normalsize\bfseries}{}{0em}{}
\titlespacing*{\section}{0pt}{2.2ex plus 1ex minus .2ex}{1.2ex plus .2ex}
\titlespacing*{\subsection}{0pt}{1.8ex plus 1ex minus .2ex}{0.9ex plus .2ex}

% --- Page furniture ----------------------------------------------------
\pagestyle{fancy}
\fancyhf{}
\renewcommand{\headrulewidth}{0.4pt}
\fancyhead[L]{\small\textcolor{headingblue}{German — Complete Notes}}
\fancyhead[R]{\small\thepage}

% --- Tables ------------------------------------------------------------
% Pandoc emits longtable; give the rows a little air.
\renewcommand{\arraystretch}{1.25}

% --- Lists -------------------------------------------------------------
\setlist[itemize]{itemsep=2pt, parsep=0pt, topsep=4pt}
\setlist[enumerate]{itemsep=2pt, parsep=0pt, topsep=4pt}

% --- Misc --------------------------------------------------------------
\setlength{\parskip}{0.6em}
\setlength{\parindent}{0pt}
\sloppy
\emergencystretch=2em
\usepackage{bookmark}
\IfFileExists{xurl.sty}{\usepackage{xurl}}{} % add URL line breaks if available
\urlstyle{same}
\hypersetup{
  pdftitle={German --- Complete Notes},
  colorlinks=true,
  linkcolor={headingblue},
  filecolor={Maroon},
  citecolor={Blue},
  urlcolor={Blue},
  pdfcreator={LaTeX via pandoc}}

\title{German --- Complete Notes}
\usepackage{etoolbox}
\makeatletter
\providecommand{\subtitle}[1]{% add subtitle to \maketitle
  \apptocmd{\@title}{\par {\large #1 \par}}{}{}
}
\makeatother
\subtitle{Everything covered so far, in study format}
\author{}
\date{Rebuilt 2026-09-18}

\begin{document}
\maketitle

{
\hypersetup{linkcolor=headingblue}
\setcounter{tocdepth}{2}
\tableofcontents
}
\section{COMPLETE NOTES --- Everything, in study
format}\label{complete-notes-everything-in-study-format}

Generated 2026-09-18 from the whole repository: GRAMMAR.md,
VOCABULARY.md, ANSWERS.md (Q1--Q23), ERROR\_LOG.md (\textasciitilde70
entries), PRONOUN\_TABLE.md, STUDY\_PLAN.md and the daily notes.

\textbf{This file is a study document, not a tracker.} It does not
replace the trackers --- it is the thing to read and revise from. If you
study only this file, you have covered everything taught so far.

\textbf{How to use it:} Part I--VII are the grammar. Part VIII is
vocabulary. Part IX is the error catalogue (your actual mistakes --- the
highest-value section). Part X is tips, traps and method.

\textbf{This file is kept live.} By standing rule, every time anything
is logged to the trackers, this document is updated too and
\texttt{build\_notes.sh} is re-run to refresh
\texttt{Complete\ Notes.tex} and \texttt{Complete\ Notes.pdf}. The PDF
is the study copy --- it should never lag behind the trackers.
\emph{(The .tex and .pdf are generated from this Markdown --- don't edit
them directly.)}

\begin{center}\rule{0.5\linewidth}{0.5pt}\end{center}

\section{PART I --- THE FOUNDATION}\label{part-i-the-foundation}

\subsection{1.1 The two core verbs}\label{the-two-core-verbs}

\textbf{sein} (to be) --- irregular, must be memorised.

{\def\LTcaptype{none} % do not increment counter
\begin{longtable}[]{@{}ll@{}}
\toprule\noalign{}
\endhead
\bottomrule\noalign{}
\endlastfoot
ich \textbf{bin} & wir \textbf{sind} \\
du \textbf{bist} & ihr \textbf{seid} \\
er/sie/es \textbf{ist} & sie/Sie \textbf{sind} \\
\end{longtable}
}

\textbf{haben} (to have) --- irregular in du/er forms.

{\def\LTcaptype{none} % do not increment counter
\begin{longtable}[]{@{}ll@{}}
\toprule\noalign{}
\endhead
\bottomrule\noalign{}
\endlastfoot
ich \textbf{habe} & wir \textbf{haben} \\
du \textbf{hast} & ihr \textbf{habt} \\
er/sie/es \textbf{hat} & sie/Sie \textbf{haben} \\
\end{longtable}
}

\textbf{!} \textbf{Your recorded mistakes here:} ``Was haben du?''
(-\textgreater{} hast), ``habst'' (-\textgreater{} hast), ``Ich haben
gegessen'' (-\textgreater{} habe), ``Er is nicht'' (-\textgreater{} ist,
English spelling leaking in), ``Du bust'' (-\textgreater{} bist).

\subsection{1.2 Feelings: two constructions, never
mixed}\label{feelings-two-constructions-never-mixed}

German expresses feelings two ways, and they cannot be combined:

\begin{itemize}
\tightlist
\item
  \textbf{sein + adjective}: Ich bin hungrig. / Ich bin müde.
\item
  \textbf{haben + noun}: Ich habe Hunger. / Ich habe Durst.
\end{itemize}

(wrong) ``Ich habe hungrig'' --- mixes the two. This was your first
logged error.

\subsection{1.3 Noun capitalization}\label{noun-capitalization}

\textbf{Every German noun is always capitalized}, everywhere in the
sentence, not just proper names: das Buch, die Katze, der Hunger, die
Zeit.

\textbf{!} Recorded slips: ``buch'', ``Hunger'' written lowercase.

But: \textbf{infinitives and separable prefixes stay lowercase} unless
genuinely nominalized. - Ich sehe \textbf{fern}. (prefix, lowercase) vs
das \textbf{Fernsehen} (the noun) - Ich kann \textbf{schwimmen}.
(infinitive) vs das \textbf{Schwimmen} (the activity)

\textbf{!} Recorded 3× as a pattern: ``Fern'', ``Schwimmen'' wrongly
capitalized.

\subsection{1.4 sie vs Sie --- the capitalization that changes
meaning}\label{sie-vs-sie-the-capitalization-that-changes-meaning}

{\def\LTcaptype{none} % do not increment counter
\begin{longtable}[]{@{}lll@{}}
\toprule\noalign{}
Form & Meaning & Verb \\
\midrule\noalign{}
\endhead
\bottomrule\noalign{}
\endlastfoot
sie (lowercase) & she & sie \textbf{ist} \\
sie (lowercase) & they & sie \textbf{sind} \\
\textbf{Sie} (always capital) & you (formal) & Sie \textbf{sind} \\
\end{longtable}
}

Mid-sentence, capitalization is the only thing distinguishing them: -
``Die Frau ist glücklich, denn \textbf{sie} hat ein Buch.'' (she) -
``\textbf{Sie} haben ein Buch.'' (you, formal)

\textbf{ihr} is never formal --- it is informal plural ``you all'' only.
Formal is always \textbf{Sie + sind}, singular or plural. \emph{(ANSWERS
Q1, Q2)}

\begin{center}\rule{0.5\linewidth}{0.5pt}\end{center}

\section{PART II --- ARTICLES, GENDER, AND THE CASE
SYSTEM}\label{part-ii-articles-gender-and-the-case-system}

\subsection{2.1 The article table --- the single most load-bearing table
in
German}\label{the-article-table-the-single-most-load-bearing-table-in-german}

{\def\LTcaptype{none} % do not increment counter
\begin{longtable}[]{@{}
  >{\raggedright\arraybackslash}p{(\linewidth - 8\tabcolsep) * \real{0.2000}}
  >{\raggedright\arraybackslash}p{(\linewidth - 8\tabcolsep) * \real{0.2000}}
  >{\raggedright\arraybackslash}p{(\linewidth - 8\tabcolsep) * \real{0.2000}}
  >{\raggedright\arraybackslash}p{(\linewidth - 8\tabcolsep) * \real{0.2000}}
  >{\raggedright\arraybackslash}p{(\linewidth - 8\tabcolsep) * \real{0.2000}}@{}}
\toprule\noalign{}
\begin{minipage}[b]{\linewidth}\raggedright
\end{minipage} & \begin{minipage}[b]{\linewidth}\raggedright
der-word (masc)
\end{minipage} & \begin{minipage}[b]{\linewidth}\raggedright
die-word (fem)
\end{minipage} & \begin{minipage}[b]{\linewidth}\raggedright
das-word (neut)
\end{minipage} & \begin{minipage}[b]{\linewidth}\raggedright
plural
\end{minipage} \\
\midrule\noalign{}
\endhead
\bottomrule\noalign{}
\endlastfoot
\textbf{Nominative} (subject) & der & die & das & die \\
\textbf{Accusative} (direct obj) & \textbf{den} & die & das & die \\
\textbf{Dative} (indirect obj) & \textbf{dem} & \textbf{der} &
\textbf{dem} & \textbf{den} (+ -n on noun) \\
\end{longtable}
}

\textbf{Notice:} in accusative, ONLY the der-word changes. That single
fact carries over everywhere --- pronouns, possessives, everything.

Indefinite: ein (der/das) / eine (die) -\textgreater{} accusative
\textbf{einen} (der-word only) -\textgreater{} dative \textbf{einem}
(der/das) / \textbf{einer} (die).

\subsection{2.2 The four cases --- what each one
marks}\label{the-four-cases-what-each-one-marks}

\begin{enumerate}
\def\labelenumi{\arabic{enumi}.}
\tightlist
\item
  \textbf{Nominative} --- the subject. Who/what is doing the verb.
\item
  \textbf{Accusative} --- the direct object. What the verb is done TO.
\item
  \textbf{Dative} --- the indirect object. The receiver/beneficiary.
\item
  \textbf{Genitive} --- possession (not yet taught).
\end{enumerate}

\textbf{!} \textbf{The misconception you formed and had corrected
(2026-09-10):} case is NOT decided by whether a verb is an ``action
verb'' or not. It is decided by the noun's ROLE in the sentence. haben,
sehen, geben --- all take accusative direct objects. Nothing about
``action-ness'' is involved.

\subsection{2.3 When accusative applies}\label{when-accusative-applies}

\textbf{Accusative = direct object of a verb.} \emph{(ANSWERS Q6)}

Verbs that take accusative objects: haben, sehen, essen, trinken,
kaufen, lesen, nehmen --- most ``doing'' verbs.

\textbf{The one big exception: sein never takes accusative.} ``To be''
doesn't act on anything, it equates two things, so nouns after sein stay
nominative.

\begin{itemize}
\tightlist
\item
  (correct) Das ist \textbf{mein} Hund. (not meinen --- sein-sentence)
\item
  (correct) Ist das \textbf{dein} Hund? (same reason)
\end{itemize}

\textbf{!} Recorded 2×: ``Ist das deinen Hund?'' --- accusative
over-applied after sein. \emph{(ANSWERS Q4)}

\textbf{Modal verbs do not change any of this.} \emph{(ANSWERS Q7)} The
second verb (the infinitive at the end) decides the case, exactly as it
would alone: - Ich möchte \textbf{den Hund} sehen. (sehen takes
accusative) - Ich kann schlafen. (schlafen takes no object)

\subsection{2.4 When dative applies --- TWO separate
causes}\label{when-dative-applies-two-separate-causes}

\textbf{Cause 1: the indirect object} (the receiver, usually a person).

\begin{quote}
\textbf{Ich gebe dem Mann ein Buch.} dem Mann = dative (receiver) · ein
Buch = accusative (thing given)
\end{quote}

\textbf{Cause 2: certain prepositions simply demand it}, regardless of
role.

\begin{quote}
\textbf{Er spielt mit mir.} --- mir isn't receiving anything; \emph{mit}
forces dative.
\end{quote}

These are genuinely two different mechanisms producing the same endings.

\subsection{\texorpdfstring{2.5 Dative-only verbs \textbf{!}
HIGH-PRIORITY WEAK
SPOT}{2.5 Dative-only verbs ! HIGH-PRIORITY WEAK SPOT}}\label{dative-only-verbs-high-priority-weak-spot}

Some verbs take a dative object even though there is only ONE object and
it looks like a direct object in English:

\begin{itemize}
\tightlist
\item
  \textbf{helfen} --- Ich helfe \textbf{ihr}. (not sie)
\item
  \textbf{danken} --- Ich danke \textbf{ihm}. (not ihn)
\item
  \textbf{gehören} --- Das Buch gehört \textbf{mir}. (not mich)
\item
  gefallen, antworten (same family)
\end{itemize}

English ``I help her'' \emph{looks} like a direct object. German treats
the person as a recipient.

\textbf{!} \textbf{This is a confirmed regression cluster.} 3/3 clean on
2026-09-10, then \textbf{1/3 on 2026-09-17} (``Ich danke ihn'' (wrong),
``Das Buch gehört mich'' (wrong)). You correctly identify the verbs as
dative-only --- you then pick the wrong form. \textbf{Drill:
mir/dir/ihm/ihr, not mich/dich/ihn/sie.}

\begin{center}\rule{0.5\linewidth}{0.5pt}\end{center}

\section{PART III --- PRONOUNS}\label{part-iii-pronouns}

\subsection{3.1 The full table}\label{the-full-table}

{\def\LTcaptype{none} % do not increment counter
\begin{longtable}[]{@{}lll@{}}
\toprule\noalign{}
Nominative & Accusative & Dative \\
\midrule\noalign{}
\endhead
\bottomrule\noalign{}
\endlastfoot
ich --- I & \textbf{mich} & \textbf{mir} \\
du --- you & \textbf{dich} & \textbf{dir} \\
er --- he & \textbf{ihn} & \textbf{ihm} \\
sie --- she & sie & \textbf{ihr} \\
es --- it & es & \textbf{ihm} \\
wir --- we & \textbf{uns} & \textbf{uns} \\
ihr --- you all & \textbf{euch} & \textbf{euch} \\
sie --- they & sie & \textbf{ihnen} \\
Sie --- you (formal) & Sie & \textbf{Ihnen} \\
\end{longtable}
}

\textbf{What to notice:} 1. Only four dative forms are genuinely new:
mir, dir, ihm, ihnen. 2. \textbf{er changes twice}
(er-\textgreater ihn-\textgreater ihm), mirroring
der-\textgreater den-\textgreater dem exactly. 3. \textbf{es and sie
barely move} --- only the dative shifts them:
es-\textgreater{}\textbf{ihm}, sie-\textgreater{}\textbf{ihr}.

\textbf{!} \textbf{es/sie dative swap --- RECURRING, 3 occurrences}
(2026-09-10, again 2026-09-17). You give \textbf{ihm} for both es and
sie. Correct: es-\textgreater ihm, sie-\textgreater{}\textbf{ihr}. This
is your most persistent pronoun error.

\subsection{3.2 The ``ihr'' trap --- three separate
jobs}\label{the-ihr-trap-three-separate-jobs}

{\def\LTcaptype{none} % do not increment counter
\begin{longtable}[]{@{}
  >{\raggedright\arraybackslash}p{(\linewidth - 4\tabcolsep) * \real{0.3333}}
  >{\raggedright\arraybackslash}p{(\linewidth - 4\tabcolsep) * \real{0.3333}}
  >{\raggedright\arraybackslash}p{(\linewidth - 4\tabcolsep) * \real{0.3333}}@{}}
\toprule\noalign{}
\begin{minipage}[b]{\linewidth}\raggedright
Form
\end{minipage} & \begin{minipage}[b]{\linewidth}\raggedright
Job
\end{minipage} & \begin{minipage}[b]{\linewidth}\raggedright
Example
\end{minipage} \\
\midrule\noalign{}
\endhead
\bottomrule\noalign{}
\endlastfoot
ihr & ``you all'' (nominative) & \textbf{Ihr} habt ein Buch \\
ihr & ``to her'' (dative) & Ich gebe \textbf{ihr} ein Buch \\
ihr/ihre & ``her'' (possessive) & \textbf{Ihre} Katze ist müde \\
\textbf{Ihr/Ihnen} (capital) & formal ``your''/``to you'' & Ist das
\textbf{Ihr} Buch? \\
\end{longtable}
}

The verb form usually disambiguates: \emph{ihr habt} (you all) vs
\emph{ihr gibt} (impossible as a subject -\textgreater{} must be
dative).

\subsection{\texorpdfstring{3.3 THE BIG PRONOUN TRAP --- gender, not
meaning \emph{(ANSWERS
Q11)}}{3.3 THE BIG PRONOUN TRAP --- gender, not meaning (ANSWERS Q11)}}\label{the-big-pronoun-trap-gender-not-meaning-answers-q11}

\textbf{A German pronoun follows the noun's GRAMMATICAL GENDER, never
the English meaning, never whether it is a person or a thing.}

Three-step method: 1. Find the noun being replaced. 2. Recall its
article --- der/die/das. 3. Convert: der-\textgreater er/ihn ·
die-\textgreater sie/sie · das-\textgreater es/es

{\def\LTcaptype{none} % do not increment counter
\begin{longtable}[]{@{}lll@{}}
\toprule\noalign{}
Sentence & Pronoun & Why \\
\midrule\noalign{}
\endhead
\bottomrule\noalign{}
\endlastfoot
Ich sehe \textbf{den Tisch} -\textgreater{} Ich sehe \textbf{ihn} & ihn
& der-word \\
Ich habe \textbf{das Buch} -\textgreater{} Ich habe \textbf{es} & es &
das-word \\
Sie kauft \textbf{die Milch} -\textgreater{} Sie kauft \textbf{sie} &
sie & die-word \\
Ich esse \textbf{einen Apfel} -\textgreater{} Ich esse \textbf{ihn} &
ihn & der Apfel \\
\end{longtable}
}

\textbf{The trap, stated plainly:} - \textbf{der Apfel} is a thing, but
takes \textbf{ihn} --- literally ``I eat him.'' Correct. - \textbf{das
Kind} is a person, but takes \textbf{es} --- ``she sees it.'' Correct. -
\textbf{das Mädchen} (girl) -\textgreater{} \textbf{es}. Yes, really.

If a German pronoun feels wrong to your English ear, that is expected
and is not a sign of error.

\textbf{Practical self-check:} say the noun WITH its article first ---
\emph{``der Apfel''} -\textgreater{} der -\textgreater{} ihn. Side
effect: if you can't recall the article, you can't choose the pronoun.
\textbf{That is why nouns must always be learned with der/die/das
attached.}

\textbf{!} Recorded 2× (same check, opposite directions): pronouns
chosen by English semantics instead of gender.

\subsection{3.4 Possessives take the same
endings}\label{possessives-take-the-same-endings}

{\def\LTcaptype{none} % do not increment counter
\begin{longtable}[]{@{}
  >{\raggedright\arraybackslash}p{(\linewidth - 6\tabcolsep) * \real{0.2500}}
  >{\raggedright\arraybackslash}p{(\linewidth - 6\tabcolsep) * \real{0.2500}}
  >{\raggedright\arraybackslash}p{(\linewidth - 6\tabcolsep) * \real{0.2500}}
  >{\raggedright\arraybackslash}p{(\linewidth - 6\tabcolsep) * \real{0.2500}}@{}}
\toprule\noalign{}
\begin{minipage}[b]{\linewidth}\raggedright
\end{minipage} & \begin{minipage}[b]{\linewidth}\raggedright
der/das-word
\end{minipage} & \begin{minipage}[b]{\linewidth}\raggedright
die-word
\end{minipage} & \begin{minipage}[b]{\linewidth}\raggedright
plural
\end{minipage} \\
\midrule\noalign{}
\endhead
\bottomrule\noalign{}
\endlastfoot
Nominative & mein Vater / \textbf{unser} Auto & mein\textbf{e} Mutter &
mein\textbf{e} Eltern \\
Accusative & mein\textbf{en} Vater & mein\textbf{e} Mutter &
mein\textbf{e} Eltern \\
Dative & mein\textbf{em} Vater & mein\textbf{er} Mutter &
mein\textbf{en} Eltern \\
\end{longtable}
}

Same pattern for dein-, sein-, ihr-, unser-, Ihr-.

\textbf{!} \textbf{Critical correction you needed twice:} the
\textbf{-en} ending is \textbf{DATIVE plural}, not ``plural in
general.'' - (correct) mit \textbf{unseren} Eltern (dative --- mit
forces it) - (correct) gegen \textbf{unsere} Eltern (accusative plural
--- NO -en)

You have twice stated the rule as ``plural always takes -en''
(2026-09-08, 2026-09-17). That is the misconception to kill.
\emph{(ANSWERS Q14)}

\subsection{\texorpdfstring{3.5 uns vs unser- \emph{(ANSWERS
Q13)}}{3.5 uns vs unser- (ANSWERS Q13)}}\label{uns-vs-unser--answers-q13}

\begin{itemize}
\tightlist
\item
  \textbf{uns} = PRONOUN, ``us'' --- Er sieht \textbf{uns}.
\item
  \textbf{unser-} = POSSESSIVE, ``our'' --- \textbf{unsere} Eltern, mit
  \textbf{unseren} Eltern
\end{itemize}

Same shape as the ihr/ihre confusion: a pronoun and its look-alike
possessive are two different words with two different jobs.

\begin{center}\rule{0.5\linewidth}{0.5pt}\end{center}

\section{PART IV --- NEGATION}\label{part-iv-negation}

\subsection{\texorpdfstring{4.1 nicht vs kein --- the decision
\emph{(ANSWERS Q3, Q8, Q9,
Q10)}}{4.1 nicht vs kein --- the decision (ANSWERS Q3, Q8, Q9, Q10)}}\label{nicht-vs-kein-the-decision-answers-q3-q8-q9-q10}

\textbf{Use kein/keine when} negating a noun that would otherwise take
ein/eine or no article (indefinite/general): - Ich habe \textbf{keinen}
Hund. (vs ``ein Hund'') - Ich habe \textbf{keine} Zeit. (Zeit normally
takes no article)

\textbf{Use nicht when:} - Negating a verb: Ich arbeite \textbf{nicht}.
- Negating an adjective: Ich bin \textbf{nicht} müde. - Negating a
\textbf{specific/definite} noun (der/die/das): Ich sehe den Hund
\textbf{nicht}.

\textbf{kein takes the same endings as ein:} keinen Hund (acc.
der-word), keine Zeit, kein Buch.

\subsection{\texorpdfstring{4.2 Why ``kein den Hund'' is impossible
\emph{(ANSWERS
Q9)}}{4.2 Why ``kein den Hund'' is impossible (ANSWERS Q9)}}\label{why-kein-den-hund-is-impossible-answers-q9}

\textbf{kein and der/die/das are mutually exclusive} --- kein REPLACES
the article slot. nicht doesn't touch the article at all.

\begin{itemize}
\tightlist
\item
  kein = article-equivalent -\textgreater{} ``Ich habe \textbf{keinen}
  Hund'' (no separate ein)
\item
  nicht = separate word, article stays -\textgreater{} ``Ich sehe
  \textbf{den} Hund \textbf{nicht}''
\end{itemize}

Two articles cannot attach to one noun, which is what ``kein den Hund''
would be.

\textbf{The real trigger:} was there originally an ein/eine (or no
article)? -\textgreater{} kein replaces it. Was there a der/die/das?
-\textgreater{} keep it, add nicht.

\subsection{\texorpdfstring{4.3 Where nicht goes \emph{(ANSWERS Q18)}
\textbf{!}
RECURRING}{4.3 Where nicht goes (ANSWERS Q18) ! RECURRING}}\label{where-nicht-goes-answers-q18-recurring}

\textbf{nicht sits as close to the end as possible, but BEFORE whatever
it negates} (verb, infinitive, or separable prefix): - Ich sehe
\textbf{nicht} fern. - Er kann \textbf{nicht} schwimmen. - Ich stehe um
sieben Uhr \textbf{nicht} auf.

\textbf{A definite object comes BEFORE nicht}, because the object is not
what is being negated --- it's established information: - (correct) Ich
sehe \textbf{den Hund nicht}. - (wrong) Ich sehe nicht den Hund. -
(correct) So kann er \textbf{den Hund nicht} sehen.

\textbf{!} You have forgotten this twice after gaps (2026-09-11,
2026-09-15) and self-corrected both times once pointed back to the rule.

\subsection{\texorpdfstring{4.4 keine Zeit vs nicht viel Zeit
\emph{(ANSWERS
Q3)}}{4.4 keine Zeit vs nicht viel Zeit (ANSWERS Q3)}}\label{keine-zeit-vs-nicht-viel-zeit-answers-q3}

Not interchangeable: - \textbf{Wir haben keine Zeit.} --- no time at
all, complete negation. - \textbf{Wir haben nicht viel Zeit.} --- some
time exists, just not much. (nicht negates \emph{viel}, the quantifier
--- consistent with the rule, not an exception.)

\begin{center}\rule{0.5\linewidth}{0.5pt}\end{center}

\section{PART V --- PREPOSITIONS}\label{part-v-prepositions}

\subsection{5.1 Dative prepositions --- always dative, no
exceptions}\label{dative-prepositions-always-dative-no-exceptions}

\textbf{mit, nach, bei, von, zu, aus, seit} (+ außer)

Mnemonic rhyme: \emph{aus, außer, bei, mit, nach, seit, von, zu}

Every one of them forces dative on whatever follows: - mit \textbf{dem}
Hund · bei \textbf{der} Arbeit · seit \textbf{dem} Jahr - von
\textbf{der} Arbeit · aus \textbf{dem} Haus · zu \textbf{meiner} Mutter

\textbf{!} \textbf{Your recurring habit: dropping the article entirely}
--- ``von Arbeit'', ``nach Wochenende'', ``von Schule''. Logged 4+
times. A dative preposition ALWAYS keeps its article.

\subsection{5.2 Accusative prepositions --- always
accusative}\label{accusative-prepositions-always-accusative}

\textbf{FUDGO: für, um, durch, gegen, ohne}

\begin{itemize}
\tightlist
\item
  Das Buch ist \textbf{für dich}.
\item
  Wir gehen \textbf{durch den} Park.
\item
  Ich esse Brot \textbf{ohne} Milch.
\item
  Wir spielen \textbf{gegen unsere} Eltern. \textbf{!} (not unseren ---
  see 3.4)
\item
  Ich stehe \textbf{um} sieben Uhr auf.
\end{itemize}

\subsection{\texorpdfstring{5.3 Two-way prepositions \textbf{!} NEEDED 4
REFRESHERS --- your hardest
topic}{5.3 Two-way prepositions ! NEEDED 4 REFRESHERS --- your hardest topic}}\label{two-way-prepositions-needed-4-refreshers-your-hardest-topic}

\textbf{in, an, auf, über, unter, vor, hinter, neben, zwischen}

\textbf{The ONLY thing to memorise --- one question:}

{\def\LTcaptype{none} % do not increment counter
\begin{longtable}[]{@{}lll@{}}
\toprule\noalign{}
Question & Meaning & Case \\
\midrule\noalign{}
\endhead
\bottomrule\noalign{}
\endlastfoot
\textbf{wohin?} (where TO?) & movement into/onto &
\textbf{ACCUSATIVE} \\
\textbf{wo?} (where?) & static location, no movement &
\textbf{DATIVE} \\
\end{longtable}
}

Worked pairs: - Ich gehe \textbf{ins} Zimmer. (motion -\textgreater{}
acc) / Ich bin \textbf{im} Zimmer. (location -\textgreater{} dat) - Ich
gehe \textbf{in den} Garten. (motion) / Der Hund ist \textbf{im} Garten.
(location) - Er legt das Buch \textbf{auf den} Tisch. (motion) / Die
Katze sitzt \textbf{auf dem} Tisch. (location)

\textbf{Contractions:} in+das -\textgreater{} \textbf{ins} · in+dem
-\textgreater{} \textbf{im} · an+das -\textgreater{} \textbf{ans} ·
an+dem -\textgreater{} \textbf{am}

\textbf{!} \textbf{Your actual friction points} (diagnosed 2026-09-15):
the wohin/wo \emph{concept} lands fine. What breaks is (a) matching the
case to the right sentence under time pressure, and (b) der-word gender
inside the phrase (der Garten -\textgreater{} dem/den, not der).

\textbf{!} \textbf{New confusion (2026-09-15):} using \textbf{nach} for
motion into a room. nach is NEVER used for entering an enclosed space.
``Ich gehe \textbf{in den} Raum,'' not ``nach der Raum.''

\subsection{\texorpdfstring{5.4 nach vs zu --- destinations
\emph{(ANSWERS Q15,
Q16)}}{5.4 nach vs zu --- destinations (ANSWERS Q15, Q16)}}\label{nach-vs-zu-destinations-answers-q15-q16}

\textbf{nach} --- destinations with NO article: - Cities/countries: nach
Berlin, nach Deutschland, nach München, nach Spanien - Directions: nach
links, nach oben - \textbf{ALSO means ``after'' (time)} --- nach dem
Wochenende, nach dem Termin. Completely different job, always dative.

\textbf{zu} --- destinations that DO take an article: - People: zu
meiner Mutter - Places: zur Schule, zum Bahnhof, zum Arzt, zur Post -
Events: zur Arbeit

\textbf{Test:} does the destination normally take der/die/das? No
-\textgreater{} nach. Yes -\textgreater{} zu.

\textbf{Contractions:} zu+dem -\textgreater{} \textbf{zum} · zu+der
-\textgreater{} \textbf{zur}

\textbf{!} \textbf{Possessives never contract} \emph{(ANSWERS Q20)}:
``zu \textbf{meiner} Mutter,'' never ``zur Mutter'' (that would mean
``to THE mother'').

\textbf{Both prepositions can appear in one sentence doing different
jobs:} \textgreater{} Wir gehen \textbf{nach dem Termin} (after ---
time) \textbf{zum Arzt} (to --- destination).

Word order: \textbf{time before place} --- see Part VI.

\textbf{Exception for later:} some countries take an article and use
``in'' --- in die Schweiz, in die USA.

\subsection{5.5 aus vs von --- both mean
``from''}\label{aus-vs-von-both-mean-from}

\begin{itemize}
\tightlist
\item
  \textbf{aus} = exiting an enclosed space, or origin/nationality

  \begin{itemize}
  \tightlist
  \item
    Ich komme \textbf{aus} Indien. · Ich komme \textbf{aus dem} Haus.
  \end{itemize}
\item
  \textbf{von} = the general ``coming from {[}a place you were{]}''

  \begin{itemize}
  \tightlist
  \item
    Sie kommt \textbf{von der} Arbeit. · Er kommt \textbf{von der}
    Schule.
  \end{itemize}
\end{itemize}

\textbf{!} Still ACTIVE, not yet drilled. ``Er kommt aus die Schule''
-\textgreater{} should be \textbf{von der} Schule.

\subsection{5.6 Fixed idioms that break the
rules}\label{fixed-idioms-that-break-the-rules}

\begin{itemize}
\tightlist
\item
  \textbf{zu Hause} = ``at home'' --- a fixed phrase, NOT zur Haus / im
  Haus / die Haus
\item
  \textbf{an der Ecke} = ``on the corner'' (location) --- an, not in
\item
  \textbf{zur Ecke} = ``to the corner'' (destination)
\end{itemize}

\begin{center}\rule{0.5\linewidth}{0.5pt}\end{center}

\section{PART VI --- WORD ORDER}\label{part-vi-word-order}

\subsection{6.1 Verb-second (V2) --- the backbone
rule}\label{verb-second-v2-the-backbone-rule}

\textbf{The conjugated verb is always in the SECOND POSITION} --- second
\emph{slot}, not second \emph{word}. Position 1 can hold a whole phrase.

{\def\LTcaptype{none} % do not increment counter
\begin{longtable}[]{@{}lll@{}}
\toprule\noalign{}
Position 1 & Position 2 (VERB) & Rest \\
\midrule\noalign{}
\endhead
\bottomrule\noalign{}
\endlastfoot
Ich & esse & heute einen Apfel \\
\textbf{Heute} & \textbf{esse} & \textbf{ich} einen Apfel \\
Am Wochenende & gehe & ich ins Kino \\
Weil ich müde bin, & bleibe & ich zu Hause \\
\end{longtable}
}

When something else takes position 1, \textbf{the subject moves after
the verb}.

\textbf{Tip:} \textbf{Your own insight (2026-09-17), and it's correct:}
fronting changes emphasis. ``Mein Vater arbeitet heute'' foregrounds the
father; ``Heute arbeitet mein Vater'' foregrounds the day.

\subsection{6.2 Yes/no questions --- verb
FIRST}\label{yesno-questions-verb-first}

Verb and subject swap places: - Du bist müde. -\textgreater{}
\textbf{Bist du} müde? - Sie wohnt in Berlin. -\textgreater{}
\textbf{Wohnt sie} in Berlin? - Er kann schwimmen. -\textgreater{}
\textbf{Kann er} schwimmen? - Du hast Zeit. -\textgreater{} \textbf{Hast
du} Zeit? \textbf{!} (not ``Habt du'')

\subsection{6.3 W-questions --- question word first, verb
second}\label{w-questions-question-word-first-verb-second}

Wer, was, wo, wann, warum, wie, wohin. \textgreater{} \textbf{Wo ist}
der Mann? · \textbf{Was hast} du? · \textbf{Wie geht's} dir?

\subsection{\texorpdfstring{6.4 Subordinate clauses --- weil sends the
verb to the END \textbf{!}
REGRESSED}{6.4 Subordinate clauses --- weil sends the verb to the END ! REGRESSED}}\label{subordinate-clauses-weil-sends-the-verb-to-the-end-regressed}

\textbf{weil} (because) pushes its own clause's verb to the very end: -
Weil ich müde \textbf{bin}, \ldots{} - Weil es \textbf{regnet}, \ldots{}

\textbf{AND the main clause still needs verb-second} --- because the
entire weil-clause occupies position 1:

\begin{quote}
\textbf{Weil es regnet, bleibe ich zu Hause.} \^{} position 1 (whole
clause) \^{} position 2 (verb) \^{} subject third
\end{quote}

(wrong) ``Weil bin ich müde, ich bleibe\ldots{}'' --- both halves wrong.

\textbf{denn} (because) is different --- it's a coordinating conjunction
like und/aber/oder, so it does NOT send the verb to the end:
\textgreater{} Ich bleibe zu Hause, \textbf{denn ich bin} müde. (normal
order both halves)

\textbf{!} \textbf{Status:} the two halves of this rule have failed in
\emph{alternation} --- 09-11 the main clause broke, 09-15 the
weil-clause broke. Both halves need to hold at once.

\subsection{\texorpdfstring{6.5 Time --- Manner --- Place \emph{(ANSWERS
Q21)}}{6.5 Time --- Manner --- Place (ANSWERS Q21)}}\label{time-manner-place-answers-q21}

After the verb, elements order as \textbf{Time -\textgreater{} Manner
-\textgreater{} Place}: \textgreater{} Ich gehe \textbf{heute} (time)
\textbf{mit meiner Mutter} (manner) \textbf{nach Berlin} (place).

Time words sit right after the conjugated verb: \textgreater{} Ich muss
\textbf{heute} schlafen.

\textbf{This is the opposite of English}, which tends to put place
before time. This is why ``Wir gehen \textbf{nach dem Termin zum Arzt}''
has the time phrase first.

\subsection{6.6 Separable verbs}\label{separable-verbs}

In a normal present-tense sentence, the \textbf{prefix splits off and
goes to the very end}: - aufstehen -\textgreater{} Ich \textbf{stehe} um
sieben Uhr \textbf{auf}. - anrufen -\textgreater{} Ich \textbf{rufe}
dich \textbf{an}. - fernsehen -\textgreater{} Ich \textbf{sehe} gern
\textbf{fern}.

\textbf{With a modal verb, it does NOT split} --- it goes to the end
whole, as an infinitive: - Ich muss um sieben Uhr \textbf{aufstehen}.

\textbf{With nicht:} nicht goes right before the prefix. - Ich stehe um
sieben Uhr \textbf{nicht auf}.

\begin{center}\rule{0.5\linewidth}{0.5pt}\end{center}

\section{PART VII --- VERBS BEYOND THE
PRESENT}\label{part-vii-verbs-beyond-the-present}

\subsection{7.1 Modal verbs}\label{modal-verbs}

{\def\LTcaptype{none} % do not increment counter
\begin{longtable}[]{@{}llll@{}}
\toprule\noalign{}
& können (can) & müssen (must) & möchten (would like) \\
\midrule\noalign{}
\endhead
\bottomrule\noalign{}
\endlastfoot
ich & kann & muss & möchte \\
du & kannst & musst & möchtest \\
er/sie/es & kann & muss & möchte \\
wir & können & müssen & möchten \\
ihr & könnt & müsst & möchtet \\
sie/Sie & können & müssen & möchten \\
\end{longtable}
}

\textbf{Structure:} the modal conjugates and sits in \textbf{position
2}; the second verb goes to the \textbf{very end in plain infinitive
form}. \textgreater{} Ich \textbf{muss} heute \textbf{schlafen}. ·
\textbf{Kann} er \textbf{schwimmen}?

Modals never have their own object --- see ANSWERS Q7 / section 2.3.

\subsection{7.2 The imperative
(commands)}\label{the-imperative-commands}

{\def\LTcaptype{none} % do not increment counter
\begin{longtable}[]{@{}
  >{\raggedright\arraybackslash}p{(\linewidth - 4\tabcolsep) * \real{0.3333}}
  >{\raggedright\arraybackslash}p{(\linewidth - 4\tabcolsep) * \real{0.3333}}
  >{\raggedright\arraybackslash}p{(\linewidth - 4\tabcolsep) * \real{0.3333}}@{}}
\toprule\noalign{}
\begin{minipage}[b]{\linewidth}\raggedright
Form
\end{minipage} & \begin{minipage}[b]{\linewidth}\raggedright
How
\end{minipage} & \begin{minipage}[b]{\linewidth}\raggedright
Example
\end{minipage} \\
\midrule\noalign{}
\endhead
\bottomrule\noalign{}
\endlastfoot
\textbf{du} & drop the -st ending & du gibst -\textgreater{}
\textbf{Gib} mir das Buch! \\
\textbf{ihr} & plain ihr-form, no pronoun & ihr esst -\textgreater{}
\textbf{Esst} das Brot! \\
\textbf{Sie} & infinitive + \textbf{Sie} stated, verb first &
\textbf{Kommen Sie} herein! \\
\end{longtable}
}

\begin{itemize}
\tightlist
\item
  Stem-change e-\textgreater i/ie \textbf{is kept}: geben
  -\textgreater{} \textbf{Gib!} · sehen -\textgreater{} \textbf{Sieh!}
\item
  Umlaut stem-changes are \textbf{dropped}: fahren -\textgreater{}
  \textbf{Fahr!} (not Fähr!)
\end{itemize}

\textbf{!} Sie-form must keep the pronoun --- ``Kommen herein'' (wrong)
-\textgreater{} ``Kommen \textbf{Sie} herein!'' (correct)

\textbf{Object cases don't change in the imperative:} ``Gib \textbf{mir}
das Buch'' --- mir is dative (receiver), das Buch accusative. Same as
any other sentence.

\subsection{\texorpdfstring{7.3 Perfekt --- the conversational past
\textbf{!} NEWEST
TOPIC}{7.3 Perfekt --- the conversational past ! NEWEST TOPIC}}\label{perfekt-the-conversational-past-newest-topic}

\textbf{Structure:} helper verb (\textbf{haben} or \textbf{sein},
conjugated, position 2) + \textbf{past participle at the very end}.

\begin{quote}
Ich \textbf{habe} gestern Fußball \textbf{gespielt}. Ich \textbf{bin}
ins Kino \textbf{gegangen}.
\end{quote}

\subsubsection{\texorpdfstring{Which helper? \emph{(ANSWERS
Q22)}}{Which helper? (ANSWERS Q22)}}\label{which-helper-answers-q22}

\textbf{sein} for: - \textbf{Motion between two places}: gehen, kommen,
fahren, fliegen, laufen - \textbf{Change of state}: aufstehen,
einschlafen, sterben, werden - \textbf{Two memorised exceptions}: sein
(ich bin gewesen), bleiben (ich bin geblieben)

\textbf{haben} for everything else --- including ``static'' verbs that
feel physical: stehen, sitzen, schlafen, spielen, arbeiten.

\textbf{Test:} moving from place A to place B, or genuinely changing
state? -\textgreater{} sein. Otherwise -\textgreater{} haben.

\textbf{!} \textbf{Counter-intuitive:} \emph{stehen} is haben (you're
not going anywhere). \emph{spielen} is haben. \emph{bleiben} is sein
despite not being motion.

\subsubsection{Participle formation}\label{participle-formation}

\begin{itemize}
\tightlist
\item
  \textbf{Regular (weak):} ge + stem + \textbf{t} -\textgreater{} machen
  -\textgreater{} ge\textbf{mach}t, spielen -\textgreater{}
  ge\textbf{spiel}t, kaufen -\textgreater{} ge\textbf{kauf}t
\item
  \textbf{Irregular (strong):} ge + (often changed) stem + \textbf{en}
  --- must be memorised per verb, like English go-\textgreater went.
\item
  \textbf{Separable verbs:} ge- goes \textbf{BETWEEN prefix and stem}
  -\textgreater{} einschlafen -\textgreater{} \textbf{ein·ge·schlafen}
  (NOT ``geeinschlafen''), aufstehen -\textgreater{}
  \textbf{auf·ge·standen}
\end{itemize}

\subsubsection{Irregular participles covered so
far}\label{irregular-participles-covered-so-far}

{\def\LTcaptype{none} % do not increment counter
\begin{longtable}[]{@{}lll@{}}
\toprule\noalign{}
Verb & Participle & Helper \\
\midrule\noalign{}
\endhead
\bottomrule\noalign{}
\endlastfoot
gehen & ge\textbf{gang}en (e-\textgreater a) & sein \\
kommen & gekommen & sein \\
fahren & ge\textbf{fahr}en (no umlaut!) & sein \\
fliegen & ge\textbf{flog}en & sein \\
bleiben & ge\textbf{blieb}en & sein \\
einschlafen & ein\textbf{ge}schlafen & sein \\
sehen & gesehen & haben \\
lesen & gelesen & haben \\
essen & gegessen & haben \\
trinken & ge\textbf{trunk}en & haben \\
stehen & ge\textbf{stand}en & haben \\
\end{longtable}
}

\textbf{!} \textbf{Your errors:} ``gefähren'' (invented umlaut
-\textgreater{} ge\textbf{fahr}en), ``gebleibt'' (regularised an
irregular verb -\textgreater{} ge\textbf{blieb}en), ``gegeinschlafen''
(ge- at the front -\textgreater{} \textbf{ein}geschlafen).

\subsubsection{\texorpdfstring{\textbf{!} THE MISCONCEPTION TO KILL
\emph{(ANSWERS
Q23)}}{! THE MISCONCEPTION TO KILL (ANSWERS Q23)}}\label{the-misconception-to-kill-answers-q23}

\textbf{sein-vs-haben and regular-vs-irregular are COMPLETELY
INDEPENDENT.} There is no correlation:

{\def\LTcaptype{none} % do not increment counter
\begin{longtable}[]{@{}lll@{}}
\toprule\noalign{}
Verb & Regular? & Helper \\
\midrule\noalign{}
\endhead
\bottomrule\noalign{}
\endlastfoot
gehen & irregular & sein \\
arbeiten & regular & \textbf{haben} \\
bleiben & irregular & sein \\
kaufen & regular & \textbf{haben} \\
kommen & irregular & sein \\
\end{longtable}
}

You proposed ``irregular verbs need sein'' --- this is false. Helper
choice is about \textbf{meaning} (motion/state-change). Participle
spelling is about \textbf{memorisation}. Two separate systems.

\subsection{7.4 Stem-changing verbs (present
tense)}\label{stem-changing-verbs-present-tense}

Some verbs change their stem vowel in \textbf{du} and \textbf{er/sie/es}
forms only:

{\def\LTcaptype{none} % do not increment counter
\begin{longtable}[]{@{}lll@{}}
\toprule\noalign{}
Verb & du & er/sie/es \\
\midrule\noalign{}
\endhead
\bottomrule\noalign{}
\endlastfoot
essen & isst & \textbf{isst} \\
sehen & siehst & \textbf{sieht} \\
lesen & liest & \textbf{liest} \\
geben & gibst & \textbf{gibt} \\
schlafen & schläfst & \textbf{schläft} \\
tragen & trägst & \textbf{trägt} \\
fahren & fährst & \textbf{fährt} \\
helfen & hilfst & \textbf{hilft} \\
\end{longtable}
}

\textbf{!} \textbf{ich never changes:} ich esse, ich sehe, ich
\textbf{helfe} (not ``hilfe''), ich gebe.

\subsection{\texorpdfstring{7.5 leben vs wohnen \emph{(ANSWERS
Q17)}}{7.5 leben vs wohnen (ANSWERS Q17)}}\label{leben-vs-wohnen-answers-q17}

\begin{itemize}
\tightlist
\item
  \textbf{wohnen} = to reside, have an address -\textgreater{} Ich
  \textbf{wohne} in Würzburg.
\item
  \textbf{leben} = to be alive / lifestyle / to live WITH people
  -\textgreater{} Er \textbf{lebt} mit seiner Mutter.
\end{itemize}

Address question -\textgreater{} wohnen. Living \emph{with} someone, or
life in general -\textgreater{} leben.

\subsection{7.6 mögen vs gern}\label{muxf6gen-vs-gern}

\begin{itemize}
\tightlist
\item
  \textbf{mögen} pairs with a \textbf{NOUN}: Ich \textbf{mag} Musik.
\item
  \textbf{gern} pairs with a \textbf{VERB}: Ich lese \textbf{gern}. /
  Ich sehe \textbf{gern} fern.
\end{itemize}

\textbf{!} Recorded 2×: ``Ich mag kochen'' -\textgreater{} should be
``Ich koche \textbf{gern}.''

\subsection{7.7 geben --- the dative/accusative double-object
verb}\label{geben-the-dativeaccusative-double-object-verb}

\begin{quote}
\textbf{Ich gebe dem Mann ein Buch.} --- dem Mann (dative receiver) +
ein Buch (accusative thing)
\end{quote}

Order: \textbf{dative before accusative} when both are nouns.

\textbf{!} \textbf{zu is redundant with geben} --- the dative already
means ``to'': (wrong) ``Er gibt das Buch \textbf{zu} seiner Mutter''
-\textgreater{} (correct) ``Er gibt \textbf{seiner Mutter} das Buch.''

\textbf{!} \textbf{geben/gehen swap} --- recorded in both directions, 2
instances.

\begin{center}\rule{0.5\linewidth}{0.5pt}\end{center}

\section{PART VIII --- VOCABULARY}\label{part-viii-vocabulary}

\subsection{8.1 Greetings \& phrases}\label{greetings-phrases}

{\def\LTcaptype{none} % do not increment counter
\begin{longtable}[]{@{}ll@{}}
\toprule\noalign{}
German & English \\
\midrule\noalign{}
\endhead
\bottomrule\noalign{}
\endlastfoot
Hallo & hello \\
Guten Morgen / Tag / Abend & good morning / day / evening \\
Gute Nacht & good night \\
Tschüss & bye (informal) \\
Auf Wiedersehen & goodbye (formal) \\
Bis später / bis bald & see you later / soon \\
Wie geht's? & how's it going? \\
Mir geht's gut & I'm doing well \\
Entschuldigung & excuse me / sorry \\
bitte & please \\
danke & thank you \\
kein Problem & no problem \\
noch mal & again / one more time \\
Gute Reise & have a good trip \\
Mach's gut & take care \\
\end{longtable}
}

\textbf{!} It is \textbf{Guten} Morgen (accusative-style set phrase),
never ``Guter Morgen.''

\subsection{8.2 People \& family}\label{people-family}

der Mann (Männer) · die Frau (Frauen) · das Kind (Kinder) · die Familie
· die Mutter (Mütter) · der Vater (Väter) · die Eltern (pl. only) · der
Bruder (Brüder) · die Schwester (Schwestern) · der Freund (Freunde)

\subsection{8.3 Food \& drink}\label{food-drink}

das Brot · das Hähnchen · der Reis · die Suppe · der Käse · die Milch ·
der Apfel (Äpfel) · das Wasser · \textbf{essen} (isst) ·
\textbf{trinken}

\textbf{!} ``Äpfel'' is the PLURAL. Singular is \textbf{Apfel} --- don't
invent umlauts.

\subsection{8.4 House \& rooms}\label{house-rooms}

das Haus (Häuser) · das Zimmer (unchanged pl.) · der Raum (Räume) · die
Küche · das Schlafzimmer · das Bad · der Tisch (Tische) · der Stuhl
(Stühle) · das Fenster (unchanged pl.)

\textbf{!} \textbf{Tisch = table, Stuhl = chair} --- swapped once.
\textbf{!} Singular is \textbf{Stuhl}, not ``Stühl.''

\textbf{Compound nouns take the gender of the LAST part:} schlafen + das
Zimmer = \textbf{das} Schlafzimmer.

\subsection{8.5 Transport, directions \& places in
town}\label{transport-directions-places-in-town}

der Bus (Busse) · die Bahn / der Zug · die Straße (Straßen) · die Ampel
(Ampeln) · die Ecke (Ecken) · \textbf{rechts} · \textbf{links} ·
\textbf{geradeaus} · die Bank (Banken) · der Park · das Rathaus · die
Post · das Kino · das Geschäft · der Bahnhof · die Schule

Phrases: \textbf{an der Ecke} (on the corner) · \textbf{zur Post} ·
\textbf{zum Bahnhof} · \textbf{ins Zentrum}

\subsection{8.6 Time \& daily routine}\label{time-daily-routine}

die Zeit · die Uhr · \textbf{heute} (today) · \textbf{morgen} (tomorrow)
· \textbf{der Morgen} (the morning --- capital M!) · das Wochenende ·
der Termin · \textbf{aufstehen} (sep.) · \textbf{arbeiten} ·
\textbf{schlafen} (schläft) · \textbf{wohnen} · \textbf{leben} ·
\textbf{um} (at, for times)

\textbf{!} morgen (tomorrow) vs \textbf{der Morgen} (the morning) ---
identical except capitalization.

\subsection{8.7 Weather}\label{weather}

das Wetter · die Sonne · der Regen · \textbf{kalt} · \textbf{warm} ·
\textbf{regnen} (impersonal: \emph{es regnet})

\subsection{8.8 Hobbies \& activities}\label{hobbies-activities}

das Hobby (Hobbys) · \textbf{lesen} (liest) · \textbf{spielen} ·
\textbf{schwimmen} · die Musik · \textbf{gern} · \textbf{fernsehen}
(sehe fern) · \textbf{mögen} (mag) · das Fernsehen (the noun)

\subsection{8.9 Clothing \& colours}\label{clothing-colours}

das T-Shirt · die Hose · das Kleid · die Jacke · der Schuh (Schuhe) ·
\textbf{tragen} (trägt) · \textbf{blau} · \textbf{rot} · \textbf{grün} ·
\textbf{schwarz}

\subsection{8.10 Shopping \& money}\label{shopping-money}

das Geld · der Preis · \textbf{kaufen} · \textbf{kosten} ·
\textbf{teuer} · \textbf{billig} · die Tasche (Taschen)

\subsection{8.11 Adjectives \& adverbs}\label{adjectives-adverbs}

müde · glücklich · hungrig · gut · schön · lang · neu · \textbf{immer}
(always) · \textbf{oft} (often) · \textbf{manchmal} (sometimes) ·
\textbf{selten} (rarely) · \textbf{nie} (never)

\subsection{8.12 Question words}\label{question-words}

wer (who) · was (what) · wo (where) · wohin (where to) · wann (when) ·
warum (why) · wie (how)

\subsection{8.13 Conjunctions}\label{conjunctions}

\textbf{und} · \textbf{aber} · \textbf{oder} · \textbf{denn} (all keep
normal word order) · \textbf{weil} (sends verb to the end)

\subsection{8.14 Numbers}\label{numbers}

eins, zwei, drei, vier, fünf, sechs, sieben, acht, neun, zehn, elf,
zwölf, dreizehn, vierzehn, fünfzehn, sechzehn, siebzehn, achtzehn,
neunzehn, zwanzig

\textbf{!} With a number, \textbf{the article is dropped}: ``zwei
Hunde,'' never ``zweinen Hunde.''

\begin{center}\rule{0.5\linewidth}{0.5pt}\end{center}

\section{PART IX --- YOUR ERROR
CATALOGUE}\label{part-ix-your-error-catalogue}

\emph{The highest-value section. These are your actual recorded
mistakes, grouped by pattern. Everything here has happened at least
once.}

\subsection{\texorpdfstring{9.1 \textbf{{[}HIGH{]}} RECURRING --- the
ones that keep coming
back}{9.1 {[}HIGH{]} RECURRING --- the ones that keep coming back}}\label{high-recurring-the-ones-that-keep-coming-back}

{\def\LTcaptype{none} % do not increment counter
\begin{longtable}[]{@{}
  >{\raggedright\arraybackslash}p{(\linewidth - 4\tabcolsep) * \real{0.3333}}
  >{\raggedright\arraybackslash}p{(\linewidth - 4\tabcolsep) * \real{0.3333}}
  >{\raggedright\arraybackslash}p{(\linewidth - 4\tabcolsep) * \real{0.3333}}@{}}
\toprule\noalign{}
\begin{minipage}[b]{\linewidth}\raggedright
Pattern
\end{minipage} & \begin{minipage}[b]{\linewidth}\raggedright
Times
\end{minipage} & \begin{minipage}[b]{\linewidth}\raggedright
The fix
\end{minipage} \\
\midrule\noalign{}
\endhead
\bottomrule\noalign{}
\endlastfoot
\textbf{3rd-person conjugated as du-form} (er hast, er siehst, er
arbeitest, er trägst, er magst) & \textbf{7+} & Load-dependent. Correct
in isolation, breaks under pressure. Embed 3rd-person subjects in mixed
practice --- do NOT drill in isolation \\
\textbf{es/sie dative swap} (sie-\textgreater ihm instead of ihr) &
\textbf{3} & es-\textgreater ihm, sie-\textgreater{}\textbf{ihr}. er and
es share ihm \\
\textbf{den used for das-words} (den Buch) & 3 & Load-dependent only;
8/8 in writing. Not a knowledge gap \\
\textbf{mein/meine gender slip} (mein Mutter) & 3 & Only in free
production. die Mutter -\textgreater{} \textbf{meine} \\
\textbf{Invented umlaut on singulars} (Bröt, Äpfel, Stühl) & 3 &
Singular is the plain form. Umlaut is the plural's job \\
\textbf{Separable prefix capitalized} (Fern, Schwimmen) & 3 & Lowercase
unless nominalized \\
\textbf{unseren used for accusative plural} & 2 & -en is \textbf{dative}
plural only \\
\textbf{Dative preposition loses its article} (von Arbeit, nach
Wochenende) & 4+ & Dative prepositions ALWAYS keep the article \\
\textbf{geben/gehen swap} & 2 & Both directions recorded \\
\textbf{mögen used with a verb} (Ich mag kochen) & 2 & mögen+noun,
gern+verb \\
\textbf{Accusative after sein} (Ist das deinen Hund?) & 2 & sein never
takes accusative \\
\end{longtable}
}

\subsection{\texorpdfstring{9.2 \textbf{{[}MED{]}} CLUSTER REGRESSIONS
--- were solid, then
broke}{9.2 {[}MED{]} CLUSTER REGRESSIONS --- were solid, then broke}}\label{med-cluster-regressions-were-solid-then-broke}

\textbf{Dative-only verbs} --- 3/3 on 09-10 -\textgreater{} \textbf{1/3
on 09-17} - (wrong) Ich danke \textbf{ihn} -\textgreater{} \textbf{ihm}
- (wrong) Das Buch gehört \textbf{mich} -\textgreater{} \textbf{mir} -
(wrong) Ich \textbf{hilfe} -\textgreater{} \textbf{helfe} (ich never
takes the stem change)

\textbf{weil word order} --- the two halves fail in alternation - 09-11:
main clause lost verb-second - 09-15: weil-clause lost verb-last

\textbf{Two-way prepositions} --- needed 4 separate refreshers (09-07,
09-08, 09-10, 09-15). Concept lands; execution under pressure breaks.

\subsection{\texorpdfstring{9.3 \textbf{{[}OK{]}} RESOLVED --- confirmed
fixed}{9.3 {[}OK{]} RESOLVED --- confirmed fixed}}\label{ok-resolved-confirmed-fixed}

\begin{itemize}
\tightlist
\item
  \textbf{geben conjugation (gibt)} --- resolved 09-07 across a genuine
  overnight gap
\item
  \textbf{verb-second general rule} --- resolved 09-03
\item
  \textbf{nicht before pronoun objects} --- ``Er sieht uns nicht''
  (correct)
\item
  \textbf{es/sie dative} --- resolved 09-10 (but recurred 09-17)
\item
  \textbf{helfen + dative} --- resolved 09-10 (but cluster recurred
  09-17)
\item
  \textbf{Sie-form imperative} --- resolved 09-15
\item
  \textbf{an der Ecke} --- resolved 09-16
\item
  \textbf{nach after-vs-to} --- resolved 09-15
\item
  \textbf{nach vs two-way in} --- resolved 09-15
\item
  \textbf{ins/im location vs motion} --- resolved 09-07
\end{itemize}

\subsection{9.4 Misconceptions you invented and had
corrected}\label{misconceptions-you-invented-and-had-corrected}

These are worth naming because inventing a plausible-but-false rule is
harder to unlearn than simply not knowing:

\begin{enumerate}
\def\labelenumi{\arabic{enumi}.}
\tightlist
\item
  (wrong) \emph{``nach is used when the noun is in infinitive form''}
  --- nouns don't have infinitives. Real rule: no-article
  -\textgreater{} nach, has-article -\textgreater{} zu.
\item
  (wrong) \emph{``Accusative is for action verbs, dative for non-action
  verbs''} --- case is decided by the noun's ROLE, not the verb's type.
\item
  (wrong) \emph{``Plural always takes -en''} --- that's \textbf{dative}
  plural specifically.
\item
  (wrong) \emph{``Irregular verbs use sein in Perfekt''} --- completely
  independent systems.
\end{enumerate}

\subsection{9.5 Anglicisms \& English
interference}\label{anglicisms-english-interference}

\begin{itemize}
\tightlist
\item
  ``Chicken'', ``Rice'' written in English -\textgreater{} das Hähnchen,
  der Reis
\item
  ``ich stay die Haus'' -\textgreater{} ich \textbf{bleibe zu Hause}
\item
  ``sehe TV'' -\textgreater{} \textbf{fernsehen}
\item
  ``Er \textbf{is} nicht'' -\textgreater{} \textbf{ist}
\item
  Writing \textbf{Munich} instead of \textbf{München}
\end{itemize}

\textbf{Tip:} \textbf{When you can't retrieve a word, write \texttt{???}
--- not the English word.} What you couldn't retrieve is the most
valuable data in the exercise.

\begin{center}\rule{0.5\linewidth}{0.5pt}\end{center}

\section{PART X --- TIPS, TRICKS \&
METHOD}\label{part-x-tips-tricks-method}

\subsection{10.1 Memory hooks}\label{memory-hooks}

{\def\LTcaptype{none} % do not increment counter
\begin{longtable}[]{@{}
  >{\raggedright\arraybackslash}p{(\linewidth - 2\tabcolsep) * \real{0.5000}}
  >{\raggedright\arraybackslash}p{(\linewidth - 2\tabcolsep) * \real{0.5000}}@{}}
\toprule\noalign{}
\begin{minipage}[b]{\linewidth}\raggedright
Hook
\end{minipage} & \begin{minipage}[b]{\linewidth}\raggedright
Covers
\end{minipage} \\
\midrule\noalign{}
\endhead
\bottomrule\noalign{}
\endlastfoot
\textbf{FUDGO} & für, um, durch, gegen, ohne (accusative preps) \\
\textbf{\emph{aus, außer, bei, mit, nach, seit, von, zu}} & dative preps
(rhyme it) \\
\textbf{wohin? / wo?} & two-way prepositions --- the ONLY thing to
memorise \\
\textbf{der -\textgreater{} den -\textgreater{} dem} mirrors \textbf{er
-\textgreater{} ihn -\textgreater{} ihm} & if one is solid, the other
follows \\
Say the noun \textbf{with its article} before choosing a pronoun & der
Apfel -\textgreater{} der -\textgreater{} ihn \\
\end{longtable}
}

\subsection{10.2 The learning rules this repo runs
on}\label{the-learning-rules-this-repo-runs-on}

\begin{itemize}
\tightlist
\item
  \textbf{Nothing is FUNCTIONAL until it survives a real overnight gap.}
  A clean first attempt is not retention --- this has been proven wrong
  twice.
\item
  \textbf{Max 2 new grammar concepts per week.} Extra hours buy depth,
  not scope.
\item
  \textbf{Every new word must appear in a sentence you produced the same
  day.}
\item
  \textbf{Retrieval beats review.} Producing from memory beats
  re-reading, always.
\item
  \textbf{Errors under load are the real data.} A fast flawed answer
  reveals more than a slow careful one.
\end{itemize}

\subsection{10.3 Writing vs voice --- they measure different
things}\label{writing-vs-voice-they-measure-different-things}

{\def\LTcaptype{none} % do not increment counter
\begin{longtable}[]{@{}lll@{}}
\toprule\noalign{}
& Writing tests & Voice tests \\
\midrule\noalign{}
\endhead
\bottomrule\noalign{}
\endlastfoot
Conjugation endings & (correct) & (wrong) STT silently repairs them \\
Article endings & (correct) & (wrong) \\
Spelling, umlauts & (correct) & (wrong) \\
Word order & (correct) & (correct) \\
\textbf{Retrieval speed} & (wrong) unlimited time & (correct) \\
\textbf{Automaticity} & (wrong) & (correct) \\
\end{longtable}
}

\textbf{Rule: learn in writing, prove in voice.} Never infer grammar
knowledge from a voice transcript --- on day one it both hid a real
error and invented a competence that wasn't there.

\subsection{10.4 Your demonstrated
strengths}\label{your-demonstrated-strengths}

Worth knowing, because they shape how to study:

\begin{enumerate}
\def\labelenumi{\arabic{enumi}.}
\tightlist
\item
  \textbf{Self-correction is your strongest skill.} Across the 09-10/11
  concepts test and the 09-15 review, you resolved most gaps yourself
  once pointed at the right reference --- that's debugging your own
  understanding, not memorising.
\item
  \textbf{You generalise patterns correctly to untaught material} ---
  produced ``Setzen Sie sich,'' ``Öffnen Sie das Fenster,'' and ``Das
  Buch gehört mir'' using verbs never formally taught.
\item
  \textbf{You notice real linguistic nuance unprompted} --- the
  Heute-fronting emphasis observation (09-17) is a genuine insight, not
  a taught rule.
\item
  \textbf{You question things that don't add up} --- asking why
  ``gegehen'' isn't the participle, catching untaught vocabulary in test
  prompts.
\end{enumerate}

\subsection{10.5 What to prioritise right
now}\label{what-to-prioritise-right-now}

\begin{enumerate}
\def\labelenumi{\arabic{enumi}.}
\tightlist
\item
  \textbf{{[}HIGH{]}} \textbf{Dative-only verbs} (helfen/danken/gehören
  + mir/dir/ihm/ihr) --- active cluster regression
\item
  \textbf{{[}HIGH{]}} \textbf{es -\textgreater{} ihm, sie
  -\textgreater{} ihr} --- 3rd recurrence
\item
  \textbf{{[}MED{]}} \textbf{weil}: verb-last inside the clause AND
  verb-second in the main clause, both at once
\item
  \textbf{{[}MED{]}} \textbf{unsere (acc. pl.) vs unseren (dat. pl.)}
  --- kill the ``-en means plural'' rule
\item
  \textbf{{[}MED{]}} \textbf{Perfekt participles} --- pure memorisation,
  no shortcut exists
\item
  \textbf{{[}OK{]}} \textbf{aus vs von} --- never yet drilled in
  isolation
\end{enumerate}

\begin{center}\rule{0.5\linewidth}{0.5pt}\end{center}

\section{APPENDIX --- QUICK REFERENCE
CARD}\label{appendix-quick-reference-card}

\textbf{Articles:} der/den/dem · die/die/der · das/das/dem ·
die/die/den+n

\textbf{Pronouns:} ich-mich-mir · du-dich-dir · er-ihn-ihm ·
sie-sie-\textbf{ihr} · es-es-\textbf{ihm} · wir-uns-uns · ihr-euch-euch
· sie-sie-ihnen

\textbf{Dative preps:} aus, außer, bei, mit, nach, seit, von, zu
\textbf{Accusative preps:} für, um, durch, gegen, ohne (FUDGO)
\textbf{Two-way:} in, an, auf, über, unter, vor, hinter, neben, zwischen
-\textgreater{} \emph{wohin?} = accusative · \emph{wo?} = dative

\textbf{Contractions:} ins (in+das) · im (in+dem) · ans · am · zum
(zu+dem) · zur (zu+der) · vom (von+dem) --- \textbf{never with
possessives}

\textbf{Word order:} verb 2nd · yes/no question verb 1st · weil verb
last · Time--Manner--Place · separable prefix to the end · nicht before
what it negates

\textbf{Perfekt:} haben/sein + participle at the end · sein =
motion/state-change + sein/bleiben · haben = everything else

\end{document}
